\section{A direct two-layer return theorem with positive physical diffusivity}
\label{sr:section}

This section proves a finite controlled return for the actual two-layer
coefficient system, including the evolving first parent. It uses the
released Alp\"oge--Buckmaster steering profile from Section 3.6.3,
pp.~18--22, and its identification in Section 3.7.6, pp.~31--32.
The estimates below are proved directly for the modified equations;
no estimate for an inviscid trajectory is transferred to a diffusive one.
The spatial realization of these amplitude equations requires the
explicit profile and localization forces calculated in the companion
sections. This section makes no assertion about an infinite sequence.

\subsection{Original parameters, entry state, and control}

Keep $A_0\geq1$, $\mu=\lambda_0=\lceil\sqrt{A_0}/2\rceil$,
and the paper's time-scaling parameter $\nu_{\rm AB}=\sqrt{A_0}$.
The physical frequencies are $K_i=\mu\widehat\lambda_i$, $i=1,2$.
The prescribed source parameters $s_*,L=L_2$ and the source profiles
$\eta,\eta_2$ are retained. Thus $0<s_*\leq1/8$,
$\eta$ is smooth nondecreasing, zero on $(-\infty,0]$ and one on
$[1,\infty)$, and $\eta_2\in C_c^\infty((0,1))$ is nonnegative with
$\int_0^1\eta_2=1$. Put $H=\|\eta_2\|_\infty$.
Physical viscosity and thermal diffusivity are both $d\geq0$.
They are distinct from $\nu_{\rm AB}$, $\mu$, and the shooting parameter
$m\in[0,3]$.

Let $t_1$ be the actual second-layer steering-entry time. Its data are
\begin{equation}\label{sr:entry}
 \begin{gathered}
 \alpha(t_1)=s_*,\quad \zeta_1(t_1)=e_2,\quad
 \Omega_1(t_1)=0,\quad \Theta_1(t_1)=-P_{1\rm in}<0,\\
 M_2(t_1)=I,\quad \zeta_2(t_1)=e(s),\quad
 \Theta_2(t_1)=-P_{2\rm in}<0,\\
 \Omega_2(t_1)=-\frac{K_2}{\sigma}P_{2\rm in}v_{\rm in}.
 \end{gathered}
\end{equation}
Here $e(s)=(\sin s,\cos s)$, $\sigma>0$ is the actual inherited scale,
$s=L\nu_{\rm AB}/\sigma$ is the actual source second angle, and
$\Gamma=\sigma\sin s$. In particular $\sigma,s,\Gamma$ are not
reset to fixed values if the preceding diffusive stage changes them.
The data $P_{1\rm in},P_{2\rm in},v_{\rm in}$ are the actual values
after growth and transition. Theorem~\ref{ave:complete_entry} proves
from both original seeds and both actual threshold crossings that these
values satisfy the explicit entry margins below, for its displayed
nonempty range of positive $d$. No coefficient comparison during
steering is assumed.

For compact notation define the following exact coordinate change:
\begin{equation}\label{sr:parameters}
 \begin{gathered}
 a=\sin s,\quad h_0=\cos s=\sqrt{1-a^2},\quad
 B_{\rm in}=\frac{K_1P_{1\rm in}}{\sigma^2},\\
 C_0=\frac{A_0\cos s_*}{\sigma^2},\quad
 C_1=\frac{A_0\sin s_*}{\sigma^2a},\\
 \delta_1=\frac{dK_1^2}{\Gamma},\qquad
 \delta_2=\frac{dK_2^2}{\Gamma},\qquad
 \tau=\Gamma(t-t_1).
 \end{gathered}
\end{equation}
Every factor used as a divisor in this change is displayed and positive; physical
time is recovered by $t=t_1+\tau/\Gamma$.
Set $T=1+L^{-1}$, $T_H=T+1/32$, and use the exact control
\begin{equation}\label{sr:z}
 z_m(\tau)=
 \begin{cases}
  1-\eta(\tau),&0\leq\tau\leq1,\\
  -mL\eta_2(L(\tau-1)),&1\leq\tau\leq T,\\
  0,&T\leq\tau\leq T_H.
 \end{cases}
 \qquad Z_m(t)=a z_m(\Gamma(t-t_1)).
\end{equation}
Flatness and interior support make these definitions smooth at both
joins. This is exactly the source pulse, with source pulse amplitude
$mL\sin s$, entry first component $a$, and entry phase length one.

\subsection{Exact closed equations and their inverse map}

Write $J=\left(\begin{smallmatrix}0&-1\\1&0\end{smallmatrix}\right)$.
The two phase gradients have the exact signed determinant
$\det(\zeta_1,\zeta_2)=-a$. Since $|\zeta_1|=1$, put
\[
 \begin{gathered}
 h=\zeta_1\cdot\zeta_2,\qquad
 D_*=h^2+a^2=|\zeta_2|^2,\qquad
 Y=\sqrt{D_*-a^2z_m^2},\\
 B=\frac{-K_1\Theta_1}{\sigma^2},\quad
 W=\frac{\Omega_1}{\sigma},\quad
 n=B+C_0+C_1h.
 \end{gathered}
\]
The symbol $D_*$ is a positive scalar, not the parent velocity matrix.
The square root selects the positive second component of $\zeta_2$.
The complete older/geometry system in $\tau$ is
\begin{equation}\label{sr:closed}
 \begin{aligned}
 h'&=W,&h(0)&=h_0,\\
 W'&=B\frac{Y-hz_m}{D_*}-\delta_1W,&W(0)&=0,\\
 B'&=-C_1W-\delta_1B,&B(0)&=B_{\rm in},\\
 n'&=-\delta_1B,&n(0)&=B_{\rm in}+C_0+C_1h_0.
 \end{aligned}
\end{equation}
The last equation is a proved consequence of the preceding ones, not
an additional independent constraint. Define $k=n/D_*$. The newest
pair is recovered from
\begin{equation}\label{sr:newest}
 \begin{gathered}
 v'=z_m-kv^2,\qquad v(0)=v_{\rm in},\\
 P_2=P_{2\rm in}\exp\!\left(\int_0^\tau(kv-\delta_2D_*)\,dr\right),\\
 \Theta_2=-P_2,\qquad
 \Omega_2=-\frac{K_2}{\sigma}P_2v.
 \end{gathered}
\end{equation}
Here $P_2$ is recovered only on an interval where $v$ has been proved
finite. The proof below constructs that interval before using this map.

For clarity, the full geometric inverse and the older physical state are
\begin{align}\label{sr:inverse}
 \zeta_2&=(az_m,Y),&
 \zeta_1&=\left(\frac{haz_m-aY}{D_*},
                    \frac{hY+a^2z_m}{D_*}\right),\\
 \phi_2&=\arcsin(az_m/\sqrt{D_*}),&
 \phi_1&=\phi_2-\arctan(a/h),&\alpha&=s_*-\phi_1,\nonumber\\
 M_1&=R_\alpha,&
 M_2&=R_{\alpha-s_*}
       \begin{pmatrix}1&-(h-h_0)/a\\0&1\end{pmatrix},&
 \Theta_1&=-\sigma^2B/K_1,&\Omega_1&=\sigma W.\nonumber
\end{align}
The expression for $M_1$ uses its original first-activation datum
$M_1=I$ when $\alpha=0$ and $p_1=e(s_*)$; it is not an artificial
reset of the first deformation at $t_1$.

\begin{proof}[Proof of the exact equations and inverse]
The actual parent is
\[
 G=-A_0R_\alpha e_2-K_1P_1\zeta_1,
 \qquad D_{<2}=\dot\alpha J+
       \Omega_1J\zeta_1\otimes\zeta_1,
 \qquad P_1=-\Theta_1.
\]
The first coefficient is $a_1=A_0\sin s_*/K_1$, because
$\phi_1+\alpha=s_*$. Therefore the diffusive first pair gives
$\dot A_1=-A_0\sin s_*\Omega_1-dK_1^2 A_1$ and
$\dot\Omega_1=-A_1\sin\phi_1-dK_1^2\Omega_1$, where
$A_1=K_1P_1$. The phase equation gives $\dot h=a\Omega_1$.
The determinant representation is $\zeta_2=h\zeta_1-aJ\zeta_1$;
solving its two components gives \eqref{sr:inverse} and
$-\sin\phi_1=a(Y-hz_m)/D_*$. Substitution and division by the
nonzero physical factors in \eqref{sr:parameters} prove the first
three equations in \eqref{sr:closed}. Differentiating
$B+C_0+C_1h$ proves the exact cancellation $n'=-\delta_1B$.

The full second coefficient retains the original base term:
\[
 a_2=\frac{aA_1+A_0(a\cos s_*+h\sin s_*)}{K_2D_*}
     =\frac{\sigma^2 a n}{K_2D_*},\qquad b_2=K_2az_m.
\]
In particular the numerator changes when the first temperature
diffuses; it has not been held fixed. The equations
$\dot\Theta_2=a_2\Omega_2-dK_2^2D_*\Theta_2$ and
$\dot\Omega_2=b_2\Theta_2-dK_2^2D_*\Omega_2$ give exactly
\eqref{sr:newest} by differentiation. Conversely differentiating its
exponential and product proves these two physical equations, so the
coordinate map and its inverse preserve the complete dynamics.

For the matrices, write $S=\left(\begin{smallmatrix}1&-(h-h_0)/a\\0&1\end{smallmatrix}\right)$.
Then $\dot S=\Omega_1(Je_2\otimes e_2)S$ because
$\dot h=a\Omega_1$. Differentiating $R_{\alpha-s_*}S$ proves
$\dot M_2=D_{<2}M_2$ and $M_2(t_1)=I$.
Furthermore $S^{-T}(a,h_0)=(a,h)$, proving
$M_2^{-T}e(s)=\zeta_2$ with its original oriented datum.
Both matrices have determinant one. Differentiating
$\zeta_{2,1}=Z_m$ in its phase equation proves the exact control
\[
 \dot\alpha=F_2^{\rm ctl}-\frac{\dot Z_m}{Y},\qquad
 F_2^{\rm ctl}=-\frac{\Omega_1\zeta_{1,1}
                       (J\zeta_1\cdot\zeta_2)}{Y}
                  =\frac{a\Omega_1\zeta_{1,1}}{Y}.
\]
Thus no extra angular control has been substituted for the source
feedback. All divisions and inverse trigonometric branches are
justified by the bounds proved next.
\end{proof}

\subsection{An explicit finite parameter range}

\begin{theorem}\label{sr:theorem}
Suppose the numerical entry and design parameters in
\eqref{sr:entry}--\eqref{sr:parameters} satisfy
\begin{equation}\label{sr:box}
 \begin{gathered}
 0<s\leq\arcsin(1/8),\quad 0<a\leq\frac18,\quad
 L\geq64,\quad 3LHa\leq\frac14,\\
 \frac12\leq B_{\rm in}\leq2,\quad
 0<C_0\leq\frac14,\quad 0<C_1\leq\frac1{32},\\
 \frac12\leq v_{\rm in}\leq2,\quad 0\leq\delta_1\leq\frac1{16}.
 \end{gathered}
\end{equation}
These are explicit inequalities on actual entry data and fixed design
numbers, not hypotheses on a subsequently unknown trajectory.
The displayed range of $s$ fixes the original positive-cosine branch;
the inequality on $\sin s$ alone would not fix that branch.
Then the entire family $m\in[0,3]$ exists through $\tau=T$,
with $P_2>0$, and its endpoint map $m\mapsto v(T;m)$ is continuous.
It satisfies the quantitative endpoint margins
\begin{equation}\label{sr:margins}
 v(T;0)\geq\frac17,\qquad v(T;3)\leq-\frac1{480}.
\end{equation}
The zero set is a nonempty compact subset of $(0,3)$ and hence
has a least element $m_*$. Every zero satisfies
\begin{equation}\label{sr:rootbounds}
 \frac{13}{96}\leq m_*\leq3-\frac1{480},\qquad
 W(T;m_*)\geq \frac{13h_0}{3840}
                     \exp(-\delta_1/L)>0.
\end{equation}
No monotonicity of the endpoint map or uniqueness of the root is claimed.

The selected state continues through the nonzero holding interval
$[T,T_H]$, with $\zeta_{2,1}=\Omega_2=0$ and with the first parent
still evolving according to \eqref{sr:closed}. On the whole selected
interval $[0,T_H]$,
\begin{equation}\label{sr:geom-bounds}
 \begin{gathered}
 h_0\leq h<3,\quad 0\leq W<9,\quad
 \frac{191}{512}<B\leq2,\quad
 \frac14<n<3,\\
 \frac{63}{64}\leq D_*<10,\quad
 Y\geq\frac{\sqrt{59}}8,\quad
 \zeta_{1,2}>\frac1{16},\quad
 \frac1{40}<k<4.
 \end{gathered}
\end{equation}
The same geometry bounds hold throughout every trial through $T$.
Every trial has $|v|<6$ through $T$, so its linear amplitude pair
never loses the negative sign of $\Theta_2$.

The exact selected logarithmic gain at every
$\tau_f\in[T,T_H]$ is
\begin{equation}\label{sr:gain}
 \log\frac{P_2(\tau_f)}{P_{2\rm in}}
 =\int_0^{T}kv\,d\tau-\delta_2\int_0^{\tau_f}D_*\,d\tau,
 \qquad
 \frac{\log3}{160}-10\delta_2\tau_f
 \leq\log\frac{P_2(\tau_f)}{P_{2\rm in}}
 \leq12T.
\end{equation}
In particular the strictly positive physical range
\begin{equation}\label{sr:physical-range}
 0<d\leq
 \min\left\{\frac{\Gamma}{16K_1^2},
 \frac{\Gamma\log3}{6400K_2^2}\right\}
\end{equation}
gives $\log(P_2(\tau_f)/P_{2\rm in})\geq\log3/320>0$.
The interval is expressed in actual physical units.
Theorem~\ref{ave:complete_entry} supplies an explicit positive
$d_{\max}$ using only the original data and chosen frequencies, and
proves \eqref{sr:box} and \eqref{sr:physical-range} for every
$0<d\leq d_{\max}$ while retaining the dependence of the actual
entry data and $\Gamma$ on $d$.
\end{theorem}

\begin{proof}
First construct geometry without using the newest amplitude pair.
On an initial interval where $B>0$, the bracket $Y-hz_m$ is
nonnegative. If $0\leq z_m\leq1$, this follows by squaring:
$Y^2-h^2z_m^2=D_*(1-z_m^2)\geq0$; if $z_m\leq0$, it follows
from $Y,h>0$. Consequently the integrating-factor formula for $W$
gives $W\geq0$. Therefore $h\geq h_0$ and
$B'=-C_1W-\delta_1B\leq0$, giving $B\leq B_{\rm in}$.
The control bound is $|az_m|\leq1/4$, so
\[
 Y^2\geq h_0^2-\frac1{16}\geq\frac{59}{64}>0.
\]
For $0\leq\tau\leq1$, $Y-hz_m\leq\sqrt{D_*}$, hence
$W'\leq B_{\rm in}/h_0$ and
$h\leq h_0+B_{\rm in}\tau^2/(2h_0)$.
On the negative pulse the same argument gives
\[
 W'\leq\frac{B_{\rm in}}{h_0}
          \{1+mL\eta_2(L(\tau-1))\}.
\]
During holding its right side is $B_{\rm in}/h_0$.
Integration on any interval up to $T_H$ therefore yields
\[
 W\leq\frac{B_{\rm in}}{h_0}(T_H+3)
 \leq\frac{26}{3}<9,
\]
because $h_0\geq15/16$, $T_H\leq17/16$, and $B_{\rm in}\leq2$.
The duration after the first part is at most
$L^{-1}+1/32\leq3/64$, so
\[
 h\leq h_0+\frac{B_{\rm in}}{2h_0}
       +\frac{B_{\rm in}}{h_0}(T_H+3)\frac3{64}
 \leq1+\frac{16}{15}+\frac{13}{32}
 =\frac{1187}{480}<3.
\]
No quotient variable has entered these estimates.

The exact integrating-factor identity for the first temperature is
\[
 B(\tau)=e^{-\delta_1\tau}B_{\rm in}
  -C_1\int_0^\tau e^{-\delta_1(\tau-r)}W(r)\,dr.
\]
Since $W\geq0$, $h'=W$, $\delta_1T_H\leq17/256$, and
$e^{-x}\geq1-x$ for $x\geq0$, it follows that
\[
 B(\tau)\geq\frac12\left(1-\frac{17}{256}\right)
          -\frac1{32}(h(\tau)-h_0)
 >\frac{239}{512}-\frac3{32}
 =\frac{191}{512}>\frac14.
\]
This strict estimate excludes a first exit through $B=0$.
All square-root denominators stay uniformly positive and $h,W,B$
remain bounded. The smooth finite-dimensional vector field is therefore
bounded with bounded first derivatives on a compact neighborhood
of every such solution. Local existence, uniqueness and the integral
equation extend the solution through $T_H$; the same neighborhood
can be chosen uniformly for $m\in[0,3]$. Differentiating that integral
equation and applying the elementary Gronwall estimate proves
continuous dependence on $m$ on this interval. This constructs the
entire older/geometry trial family before the newest pair.

Now $n=B+C_0+C_1h>1/4$, and
$n'=-\delta_1B\leq0$, so
$n\leq B_{\rm in}+C_0+C_1h_0\leq73/32<3$.
The asserted $D_*$ and $k$ bounds follow. For the first phase,
$|az_m|\leq1/4$, $a\leq1/8$, $h\geq15/16$ and $Y>7/8$
give
\[
 \zeta_{1,2}=\frac{hY+a^2z_m}{D_*}
 >\frac{(15/16)(7/8)-1/32}{10}
 =\frac{101}{1280}>\frac1{16}.
\]
Thus the branches and control denominators in \eqref{sr:inverse}
stay in their specified domains. This proves all geometry claims.

Construct the Riccati variable next. On $[0,1]$, its forcing
$z_m=1-\eta$ is nonnegative. The reciprocal differential inequality
$(1/v)'\leq4$, initially while $v>0$, gives
\begin{equation}\label{sr:v-ramp}
 \frac{v_{\rm in}}{1+4v_{\rm in}\tau}
 \leq v(\tau)\leq v_{\rm in}+\tau\leq3.
\end{equation}
The positive lower bound rules out a first zero, and the upper bound
follows by integrating $v'\leq1$. These bounds also exclude finite
escape. Put $v_d=v(1)$. In particular $v_d\geq1/6$.
Moreover
\[
 \int_0^1kv^2\,d\tau
 \geq\frac1{40}\int_0^1
       \frac{v_{\rm in}^2}{(1+4v_{\rm in}\tau)^2}\,d\tau
 =\frac{v_{\rm in}^2}{40(1+4v_{\rm in})}
 \geq\frac1{480}.
\]
The last inequality follows because $x^2/(1+4x)$ is increasing
for $x>0$, as its derivative is $(2x+4x^2)/(1+4x)^2>0$.
Integrating the Riccati equation proves
\begin{equation}\label{sr:vd}
 \frac16\leq v_d\leq3-\frac1{480}.
\end{equation}

On the pulse put $V(y)=v(1+y/L)$ and
$M_2(y)=\int_0^y\eta_2(r)\,dr$. Its exact integral equation is
\[
 V(y)=v_d-mM_2(y)
       -\frac1L\int_0^y k(1+r/L)V(r)^2\,dr.
\]
Under the temporary bound $|V|\leq6$, it gives
$V\leq v_d<3$ and
$V\geq-3-144/L\geq-21/4>-6$.
The bound is strict initially and strictly improves itself at every
possible first exit, so $|V|<6$ on the full pulse. The quotient has
therefore been constructed with a uniform finite bound, and the
exponential formula \eqref{sr:newest} constructs a negative
$\Theta_2$ solving the entire original linear pair. Its uniqueness
follows by subtracting two solutions of that linear pair and applying
the integral Gronwall inequality to their zero initial difference.
In particular the construction agrees with the unique linear-pair
solution and has not assumed absence of a temperature zero.

For $m=0$, the pulse solves $v'=-kv^2$, so
$v(T;0)\geq v_d/(1+4v_d/L)\geq1/(6+4/L)>1/7$.
For $m=3$, the integral equation and \eqref{sr:vd} give
$v(T;3)\leq v_d-3\leq-1/480$.
The uniformly finite quotient and geometry imply continuous
dependence of $v(T;m)$ on $m$ by their integral equations.
The intermediate value theorem proves existence of a zero. Its
zero set is closed in $[0,3]$ and excludes both endpoints, proving
the stated compactness and least-root selection.

At any selected zero, $v'=-mL\eta_2-kv^2\leq0$ throughout
the pulse. Thus $0\leq v\leq v_d$ there, and
\[
 m=v_d-\int_1^Tkv^2\,d\tau
 \in\left[v_d-\frac{4v_d^2}{L},v_d\right]
 \subseteq\left[\frac{13}{96},3-\frac1{480}\right].
\]
For the lower inclusion use $v_d\leq3$, $L\geq64$ and
$v_d\geq1/6$. On the pulse,
\[
 W'+\delta_1W
 =B\frac{Y+hmL\eta_2}{D_*}
 \geq\frac{h_0}{40}mL\eta_2(L(\tau-1)).
\]
Its integrating-factor formula, $W(1)\geq0$ and the unit integral
of $\eta_2$ prove \eqref{sr:rootbounds}.
On holding, $z_m=0$ and $v(T)=0$, so uniqueness gives $v=0$.
The older geometry continues by the already established bounds,
and \eqref{sr:newest} gives the exact remaining temperature decay.

Finally, \eqref{sr:v-ramp} gives
\[
 \int_0^1kv\,d\tau\geq
 \frac1{40}\int_0^1
 \frac{v_{\rm in}}{1+4v_{\rm in}\tau}\,d\tau
 =\frac1{160}\log(1+4v_{\rm in})
 \geq\frac{\log3}{160}.
\]
The selected pulse adds nonnegative gain, the hold adds zero gain,
and $v\leq3$, $k<4$ bound the gain above by $12T$.
The exact diffusion exponent is at most $10\delta_2\tau_f$.
This proves \eqref{sr:gain}; substitution of
\eqref{sr:physical-range} gives its positive lower bound.
\end{proof}

\subsection{Returned state, holding dynamics, and diffusion costs}

At the selected physical return time $t_b=t_1+T/\Gamma$ one has
\[
 \zeta_2(t_b)=(0,\sqrt{D_*(T)}),\qquad\Omega_2(t_b)=0,
 \quad\zeta_1(t_b)=\frac{(-a,h(T))}{\sqrt{D_*(T)}},\qquad
 \alpha(t_b)=s_*+\arctan(a/h(T)),
\]
and $\Theta_1(t_b)=-\sigma^2B(T)/K_1<0$,
$\Omega_1(t_b)=\sigma W(T)>0$. The first parent is therefore
neither frozen nor returned to zero. Its full motion creates the
specified returned state of the second wave.
The controlled holding interval has physical duration $1/(32\Gamma)$.
On it the complete system is
\[
 \begin{gathered}
 h'=W,\quad W'=B/\sqrt{h^2+a^2}-\delta_1W,
 \quad B'=-C_1W-\delta_1B,\\
 P_2(\tau)=P_2(T)\exp\!\left(-\delta_2
                  \int_T^\tau(h(r)^2+a^2)\,dr\right).
 \end{gathered}
\]
These formulas retain the evolving phase length inside the diffusion
integral. They also prove the exact stationarity of
$\zeta_{2,1},\Omega_2$ throughout this finite hold.

The returned parent vorticity also gives positive scalar curl of the
full local velocity at the origin. Indeed $z_m=z_m'=0$ on holding,
so the proved control and inverse give the exact physical formulas
\begin{equation}\label{sr:origin_vorticity}
 \dot\alpha=-\frac{a^2\Omega_1}{h^2+a^2},\qquad
 \omega(0,t)=2\dot\alpha+\Omega_1
   =\Omega_1\left(1-\frac{2a^2}{h^2+a^2}\right)
   \geq\frac{31}{32}\Omega_1>0.
\end{equation}
For the inequality use $h\geq h_0$, $h_0^2+a^2=1$ and
$a^2\leq1/64$. The positive forcing in the holding equation for
$W$ gives $W(\tau)\geq W(T)e^{-\delta_1(\tau-T)}>0$ by its
integrating factor. These statements include the return endpoint
because the original pulse is identically zero near that endpoint.
The compact spatial construction has constant envelopes and exact
linear profiles near the origin, so its curl there is precisely the
quantity in \eqref{sr:origin_vorticity}; no cutoff derivative enters it.

The total first-wave damping exponents on the ramp, pulse and hold are
$\delta_1$, $\delta_1/L$ and $\delta_1/32$, respectively.
The corresponding second-wave exponents are exactly
\[
 \delta_2\int_0^1D_*\,d\tau,\qquad
 \delta_2\int_1^T D_*\,d\tau,\qquad
 \delta_2\int_T^{T_H}D_*\,d\tau.
\]
Each exponent lies between $(63/64)\delta_2$ and $10\delta_2$
times its displayed interval length. In original physical coordinates
these equal $dK_2^2\int|\zeta_2(t)|^2dt$, with
$K_2=\mu\widehat\lambda_2$ and $dt=d\tau/\Gamma$; none of
the factors $\mu$, $\nu_{\rm AB}$, $K_i$ or $\Gamma$ has been dropped.
The exact first-amplitude identity in the proof also records the
additional parent-response term
$-C_1\int e^{-\delta_1(\tau-r)}W(r)dr$; a pure common-damping
multiplier would miss that term.

For the parent affine fields the altered gradient is not independently
prescribed. With $B_1=-\sigma^2B\zeta_1$, direct differentiation gives
\[
 \dot G+D_{<2}^{T}G=-dK_1^2B_1,
 \qquad C=2\dot\alpha+\Omega_1,
 \qquad\dot C-G_1=2\ddot\alpha-A_0\sin\alpha-dK_1^2\Omega_1.
\]
Thus $G\cdot x,D_{<2}x$ have scalar force
$(-dK_1^2B_1)\cdot x$ and full vector force
$\tfrac12(2\ddot\alpha-A_0\sin\alpha-dK_1^2\Omega_1)Jx$
with pressure
$-\tfrac12x^T\operatorname{sym}(\dot D_{<2}+D_{<2}^2-e_2\otimes G)x$.
Their Laplacians vanish. These affine formulas are local exact fields;
the compact original-profile realization and its additional spatial
forces are supplied by the accompanying finite-profile construction.
They must not be discarded when interpreting the amplitude result
as a statement about the full spatial equations.

Here is a direct verification of the affine force and pressure, including
their signs. Put $D=D_{<2}$, $u(x,t)=D(t)x$ and
$\theta(x,t)=G(t)\cdot x$. Both summands in $D$ have zero trace, so
$\operatorname{div}u=0$ and
$D^2=-\det(D)I$ by multiplication of a general trace-free
$2\times2$ matrix. The scalar residual is exactly
$(\dot G+D^TG)\cdot x$. In that derivative the term
$K_1\dot\Theta_1\zeta_1$ cancels
$(\Omega_1J\zeta_1\otimes\zeta_1)^T(-A_0R_\alpha e_2)$,
apart from $-dK_1^2B_1$; the angular terms cancel because
$\dot\zeta_1=\dot\alpha J\zeta_1$ and
$\partial_t(R_\alpha e_2)=\dot\alpha J R_\alpha e_2$.
Also $\zeta_1\cdot J\zeta_1=0$ eliminates the self-interaction.
For momentum let
$\mathcal A=\dot D+D^2-e_2\otimes G$.
Then
$\partial_tu+(u\cdot\nabla)u-\theta e_2=\mathcal A x$,
whereas the displayed pressure contributes
$-\operatorname{sym}(\mathcal A)x$.
Their sum is $\operatorname{skew}(\mathcal A)x$.
The scalar curl of $Dx$ is $C=D_{21}-D_{12}$; since $D^2$ is
scalar and the curl of $(G\cdot x)e_2$ is $G_1$,
\[
 \operatorname{skew}(\mathcal A)
 =\frac{\dot C-G_1}{2}J.
\]
Finally $\dot\Omega_1=-K_1P_1\zeta_{1,1}-dK_1^2\Omega_1$
and $G_1=A_0\sin\alpha-K_1P_1\zeta_{1,1}$ give precisely
the force stated above. Both affine Laplacians are zero by direct
second differentiation. Thus these are the full vector equations with
the specified pressure, not merely a vorticity residual.
